% GENERATED FILE. Edit canonical lessons, then run npm run generate:books.
% canonical-source-hash: fnv1a64:e62eb6a378d4da80
% canonical-lessons: FR-C523-episode, FR-C523-saison, FR-C523-suspense, FR-C523-plateforme, FR-C523-palpitant

\chapter{Une série à suivre}
\label{ch:fr-une-serie-a-suivre}
\hlchaptermodality{523}

\begin{chapteropening}
\textbf{By the end of this chapter:} \emph{I can talk about a series I am watching.}

Complete the last lesson of chapter 523: I can talk about a series I am watching.
\end{chapteropening}

\section[l'épisode]{l'épisode — episode}

\label{lesson:FR-C523-episode}



\begin{quote}
Before the new one: say the French for biography, then the French for tale.
\end{quote}

\subsection*{You'll want to know: l'épisode}
\textbf{l'épisode} — \textquotedblleft{}episode\textquotedblright{}.

\begin{quote}
J'ai regardé trois épisodes hier soir.
\end{quote}

\begin{quote}
\emph{I watched three episodes last night.}
\end{quote}

Masculine.

\begin{cousinweb}
\emph{l'épisode} comes from Greek \emph{epeisódion} \textquotedblleft{}an added part\textquotedblright{} (EtymDB), built on \emph{epí} \textquotedblleft{}upon\textquotedblright{} and \emph{eísodos} \textquotedblleft{}an entrance\textquotedblright{}.
\end{cousinweb}

\begin{grammarlens}[title={What you've built}]
One more word for telling, explaining and describing at B1.
\end{grammarlens}

\subsection*{Your turn}
\begin{itemize}
\raggedright
  \item \emph{Say it:} \emph{l'épisode}
  \item \emph{Say it:} \emph{l'épisode}, once more
  \item \emph{Say it:} a sentence of your own with \emph{l'épisode}
  \item \emph{Recall:} say the French for biography, then the French for tale, then say \emph{l'épisode} again
\end{itemize}

\begin{tcolorbox}[breakable,colback=teal!4,colframe=teal!35!black,title={Before you move on}]
What does l'épisode mean? (\textquotedblleft{}Episode\textquotedblright{}.) Say it once more.
\end{tcolorbox}
\section[la saison]{la saison — season}

\label{lesson:FR-C523-saison}



\begin{quote}
Before the new one: say the French for tale, then the French for episode.
\end{quote}

\subsection*{You'll want to know: la saison}
\textbf{la saison} — \textquotedblleft{}season\textquotedblright{}.

\begin{quote}
La deuxième saison de la série sort en mars.
\end{quote}

\begin{quote}
\emph{The second season of the series comes out in March.}
\end{quote}

Also the seasons of the year, and \emph{la haute saison} for tourism.

\begin{cousinweb}
\emph{la saison} is inherited from Latin \emph{satiō}, \emph{satiōnem} \textquotedblleft{}a sowing\textquotedblright{} (EtymDB), from \emph{serere} \textquotedblleft{}to sow\textquotedblright{} (Lewis \& Short): first the time for sowing.
\end{cousinweb}

\begin{grammarlens}[title={What you've built}]
One more word for telling, explaining and describing at B1.
\end{grammarlens}

\subsection*{Your turn}
\begin{itemize}
\raggedright
  \item \emph{Say it:} \emph{la saison}
  \item \emph{Say it:} \emph{la saison}, once more
  \item \emph{Say it:} a sentence of your own with \emph{la saison}
  \item \emph{Recall:} say the French for tale, then the French for episode, then say \emph{la saison} again
\end{itemize}

\begin{tcolorbox}[breakable,colback=teal!4,colframe=teal!35!black,title={Before you move on}]
What does la saison mean? (\textquotedblleft{}Season\textquotedblright{}.) Say it once more.
\end{tcolorbox}
\section[le suspense]{le suspense — suspense}

\label{lesson:FR-C523-suspense}



\begin{quote}
Before the new one: say the French for episode, then the French for season.
\end{quote}

\subsection*{You'll want to know: le suspense}
\textbf{le suspense} — \textquotedblleft{}suspense\textquotedblright{}.

\begin{quote}
Ce film garde le suspense jusqu'à la fin.
\end{quote}

\begin{quote}
\emph{This film keeps up the suspense to the end.}
\end{quote}

Said the English way or \emph{sus-pense}. A different word from \emph{en suspens} \textquotedblleft{}pending\textquotedblright{}.

\begin{cousinweb}
\emph{le suspense} was borrowed from English \emph{suspense}, which goes back to Latin \emph{suspendere} \textquotedblleft{}to hang up\textquotedblright{}, built on \emph{pendere} \textquotedblleft{}to hang\textquotedblright{} (Lewis \& Short).
\end{cousinweb}

\begin{grammarlens}[title={What you've built}]
One more word for telling, explaining and describing at B1.
\end{grammarlens}

\subsection*{Your turn}
\begin{itemize}
\raggedright
  \item \emph{Say it:} \emph{le suspense}
  \item \emph{Say it:} \emph{le suspense}, once more
  \item \emph{Say it:} a sentence of your own with \emph{le suspense}
  \item \emph{Recall:} say the French for episode, then the French for season, then say \emph{le suspense} again
\end{itemize}

\begin{tcolorbox}[breakable,colback=teal!4,colframe=teal!35!black,title={Before you move on}]
What does le suspense mean? (\textquotedblleft{}Suspense\textquotedblright{}.) Say it once more.
\end{tcolorbox}
\section[la plateforme]{la plateforme — (streaming) platform; platform}

\label{lesson:FR-C523-plateforme}



\begin{quote}
Before the new one: say the French for season, then the French for suspense.
\end{quote}

\subsection*{You'll want to know: la plateforme}
\textbf{la plateforme} — \textquotedblleft{}(streaming) platform; platform\textquotedblright{}.

\begin{quote}
Cette série est disponible sur une plateforme.
\end{quote}

\begin{quote}
\emph{This series is available on a streaming platform.}
\end{quote}

Also a station platform in some contexts, but that is usually \emph{le quai}. Older spelling \emph{plate-forme}.

\begin{cousinweb}
\emph{la plateforme} is built from \emph{plat} \textquotedblleft{}flat\textquotedblright{} and \emph{forme} \textquotedblleft{}shape\textquotedblright{}, from Latin \emph{forma} (Lewis \& Short).
\end{cousinweb}

\begin{grammarlens}[title={What you've built}]
One more word for telling, explaining and describing at B1.
\end{grammarlens}

\subsection*{Your turn}
\begin{itemize}
\raggedright
  \item \emph{Say it:} \emph{la plateforme}
  \item \emph{Say it:} \emph{la plateforme}, once more
  \item \emph{Say it:} a sentence of your own with \emph{la plateforme}
  \item \emph{Recall:} say the French for season, then the French for suspense, then say \emph{la plateforme} again
\end{itemize}

\begin{tcolorbox}[breakable,colback=teal!4,colframe=teal!35!black,title={Before you move on}]
What does la plateforme mean? (\textquotedblleft{}(streaming) platform; platform\textquotedblright{}.) Say it once more.
\end{tcolorbox}
\section[palpitant]{palpitant — thrilling, exciting}

\label{lesson:FR-C523-palpitant}



\begin{quote}
Before the new one: say the French for suspense, then the French for (streaming) platform.
\end{quote}

\subsection*{You'll want to know: palpitant}
\textbf{palpitant} — \textquotedblleft{}thrilling, exciting\textquotedblright{}.

\begin{quote}
Le dernier épisode était palpitant.
\end{quote}

\begin{quote}
\emph{The last episode was thrilling.}
\end{quote}

Built on \emph{palpiter} \textquotedblleft{}to throb\textquotedblright{}.

\begin{cousinweb}
\emph{palpitant} is the present participle of \emph{palpiter} \textquotedblleft{}to throb\textquotedblright{}, a learned borrowing from Latin \emph{palpitāre} \textquotedblleft{}to flutter, to throb\textquotedblright{}, from \emph{palpāre} \textquotedblleft{}to touch softly\textquotedblright{} (Lewis \& Short).
\end{cousinweb}

\begin{grammarlens}[title={What you've built}]
One more word for telling, explaining and describing at B1.
\end{grammarlens}

\subsection*{Your turn}
\begin{itemize}
\raggedright
  \item \emph{Say it:} \emph{palpitant}
  \item \emph{Say it:} \emph{palpitant}, once more
  \item \emph{Say it:} a sentence of your own with \emph{palpitant}
  \item \emph{Recall:} say the French for suspense, then the French for (streaming) platform, then say \emph{palpitant} again
\end{itemize}

\begin{tcolorbox}[breakable,colback=teal!4,colframe=teal!35!black,title={Before you move on}]
What does palpitant mean? (\textquotedblleft{}Thrilling, exciting\textquotedblright{}.) Say it once more.
\end{tcolorbox}
